\chapter{Resultados del estudio anidado}\label{sec:resultados-experimentales}

\section{Resultados principales: desempeño y descriptor individual}

La Tabla~\ref{tab:resultados-confirmatorios} compara las cinco estrategias con la biblioteca de 21 descriptores. Individual corresponde a una política de selección interna: el descriptor elegido puede cambiar entre particiones. El panel B identifica esos descriptores y sus frecuencias, evitando atribuir la media a un único extractor cuando hubo varios. La Tabla~\ref{tab:accuracy} presenta la exactitud bajo el mismo protocolo.

\begin{resulttable}
\centering\fontsize{9}{11}\selectfont\setlength{\tabcolsep}{3.6pt}
\caption{Macro-F1 externo de las cinco estrategias con la biblioteca de 21 descriptores. Media y desviación estándar entre particiones; Outex tiene un único split. El máximo de cada fila se destaca sin implicar significancia estadística.}
\label{tab:resultados-confirmatorios}
\textbf{Panel A: macro-F1 externo}\par\smallskip
\begin{tabular}{@{}llccccc@{}}
\toprule
Dataset & Clasif. & Individual & Completa & Top-$k$ & GFS & Homogénea \\
\midrule
DTD & SVM & $0{,}8426\,\pm\,0{,}0109$ & $0{,}8629\,\pm\,0{,}0080$ & $0{,}8684\,\pm\,0{,}0080$ & $\mathbf{0{,}8694}\,\pm\,0{,}0096$ & $0{,}8606\,\pm\,0{,}0086$ \\
DTD & ResMLP & $0{,}8373\,\pm\,0{,}0089$ & $0{,}8608\,\pm\,0{,}0100$ & $0{,}8648\,\pm\,0{,}0084$ & $\mathbf{0{,}8660}\,\pm\,0{,}0095$ & $0{,}8611\,\pm\,0{,}0060$ \\
FMD & SVM & $0{,}9726\,\pm\,0{,}0150$ & $0{,}9644\,\pm\,0{,}0174$ & $\mathbf{0{,}9820}\,\pm\,0{,}0105$ & $0{,}9766\,\pm\,0{,}0135$ & $0{,}9789\,\pm\,0{,}0122$ \\
FMD & ResMLP & $0{,}9690\,\pm\,0{,}0133$ & $0{,}9518\,\pm\,0{,}0203$ & $0{,}9753\,\pm\,0{,}0140$ & $\mathbf{0{,}9783}\,\pm\,0{,}0123$ & $0{,}9773\,\pm\,0{,}0109$ \\
CUReT & SVM & $0{,}9948\,\pm\,0{,}0028$ & $\mathbf{0{,}9989}\,\pm\,0{,}0015$ & $0{,}9989\,\pm\,0{,}0010$ & $0{,}9988\,\pm\,0{,}0008$ & $0{,}9989\,\pm\,0{,}0010$ \\
CUReT & ResMLP & $0{,}9971\,\pm\,0{,}0030$ & $\mathbf{0{,}9993}\,\pm\,0{,}0010$ & $0{,}9986\,\pm\,0{,}0015$ & $0{,}9980\,\pm\,0{,}0013$ & $0{,}9982\,\pm\,0{,}0015$ \\
Outex & SVM & $0{,}9066$ & $\mathbf{0{,}9674}$ & $0{,}9441$ & $0{,}9610$ & $0{,}9307$ \\
Outex & ResMLP & $0{,}9476$ & $0{,}9537$ & $0{,}9607$ & $\mathbf{0{,}9715}$ & $0{,}9306$ \\
\bottomrule
\end{tabular}

\par\medskip\textbf{Panel B: descriptor Individual seleccionado}\par\smallskip
\label{tab:individual-identidad}
\begin{tabular}{@{}lll@{}}
\toprule
Dataset & Clasificador & Descriptor (frecuencia) \\
\midrule
DTD & SVM & SigLIP-B (8/10); DINOv2-L (2/10) \\
DTD & ResMLP & SigLIP-B (10/10) \\
FMD & SVM & SigLIP-B (15/15) \\
FMD & ResMLP & SigLIP-B (15/15) \\
CUReT & SVM & DINOv2-L (2/2) \\
CUReT & ResMLP & DINOv2-S (1/2); DINOv2-L (1/2) \\
Outex & SVM & DINOv2-L (1/1) \\
Outex & ResMLP & N-gramas + SVD (1/1) \\
\bottomrule
\end{tabular}
\par\smallskip{\footnotesize Frecuencias de selección interna; cuando cambia el descriptor, la media del panel A reúne esas elecciones.}
\end{resulttable}
\FloatBarrier

\begin{resulttable}
\centering\fontsize{8}{10}\selectfont\setlength{\tabcolsep}{3.0pt}
\caption{Exactitud externa de cinco estrategias con la biblioteca común de 22 descriptores, incluidos RGB N-gramas + SVD y BEiTv2-B. Se muestran medias entre condiciones externas; Outex tiene un único split oficial.}
\label{tab:accuracy}
\begin{tabular}{@{}llccccc@{}}
\toprule
Dataset & Clasif. & Individual & Completa & Top-$k$ & GFS & Homogénea \\
\midrule
DTD & SVM & $0{,}8448$ & $0{,}8654$ & $0{,}8714$ & $\mathbf{0{,}8716}$ & $0{,}8621$ \\
DTD & ResMLP & $0{,}8382$ & $0{,}8635$ & $\mathbf{0{,}8652}$ & $0{,}8651$ & $0{,}8626$ \\
FMD & SVM & $0{,}9727$ & $0{,}9647$ & $\mathbf{0{,}9820}$ & $0{,}9767$ & $0{,}9790$ \\
FMD & ResMLP & $0{,}9690$ & $0{,}9563$ & $0{,}9740$ & $\mathbf{0{,}9783}$ & $0{,}9773$ \\
CUReT & SVM & $0{,}9948$ & $\mathbf{0{,}9989}$ & $\mathbf{0{,}9989}$ & $0{,}9988$ & $\mathbf{0{,}9989}$ \\
CUReT & ResMLP & $0{,}9971$ & $\mathbf{0{,}9991}$ & $0{,}9984$ & $0{,}9986$ & $0{,}9982$ \\
Outex & SVM & $0{,}9103$ & $\mathbf{0{,}9662}$ & $0{,}9441$ & $0{,}9603$ & $0{,}9324$ \\
Outex & ResMLP & $0{,}9485$ & $0{,}9456$ & $0{,}9662$ & $\mathbf{0{,}9721}$ & $0{,}9338$ \\
Soil & SVM & $0{,}9044$ & $\mathbf{0{,}9281}$ & $0{,}9225$ & $0{,}9152$ & $0{,}9126$ \\
Soil & ResMLP & $0{,}9044$ & $\mathbf{0{,}9184}$ & $0{,}9181$ & $0{,}9170$ & $0{,}9135$ \\
KTH-TIPS2-b & SVM & $0{,}9415$ & $\mathbf{0{,}9609}$ & $0{,}9419$ & $0{,}9112$ & $0{,}9293$ \\
KTH-TIPS2-b & ResMLP & $0{,}9099$ & $\mathbf{0{,}9499}$ & $0{,}9346$ & $0{,}9219$ & $0{,}9165$ \\
\bottomrule
\end{tabular}
\end{resulttable}
\FloatBarrier

\begin{resulttable}
\centering\fontsize{8}{10}\selectfont\setlength{\tabcolsep}{3.0pt}
\caption{Dimensión media de las representaciones usadas en la evaluación externa con 22 descriptores.}
\label{tab:dimensiones}
\begin{tabular}{@{}llrrrrr@{}}
\toprule
Dataset & Clasif. & Individual & Completa & Top-$k$ & GFS & Homogénea \\
\midrule
DTD & SVM & $819{,}2$ & $20652{,}0$ & $4480{,}0$ & $5390{,}4$ & $2496{,}0$ \\
DTD & ResMLP & $768{,}0$ & $20652{,}0$ & $4211{,}2$ & $4626{,}4$ & $2598{,}4$ \\
FMD & SVM & $768{,}0$ & $20652{,}0$ & $2969{,}6$ & $2687{,}5$ & $2338{,}1$ \\
FMD & ResMLP & $768{,}0$ & $20652{,}0$ & $3404{,}8$ & $3077{,}5$ & $2218{,}7$ \\
CUReT & SVM & $1024{,}0$ & $20652{,}0$ & $3328{,}0$ & $4402{,}0$ & $2944{,}0$ \\
CUReT & ResMLP & $704{,}0$ & $20652{,}0$ & $2176{,}0$ & $2880{,}0$ & $1856{,}0$ \\
Outex & SVM & $1024{,}0$ & $20652{,}0$ & $7808{,}0$ & $6784{,}0$ & $3712{,}0$ \\
Outex & ResMLP & $256{,}0$ & $20652{,}0$ & $6016{,}0$ & $3346{,}0$ & $2176{,}0$ \\
Soil & SVM & $768{,}0$ & $20652{,}0$ & $4898{,}1$ & $3709{,}9$ & $2184{,}5$ \\
Soil & ResMLP & $768{,}0$ & $20652{,}0$ & $5529{,}6$ & $4415{,}6$ & $2235{,}7$ \\
KTH-TIPS2-b & SVM & $832{,}0$ & $20652{,}0$ & $1280{,}0$ & $960{,}0$ & $960{,}0$ \\
KTH-TIPS2-b & ResMLP & $672{,}0$ & $20652{,}0$ & $1600{,}0$ & $1440{,}0$ & $928{,}0$ \\
\bottomrule
\end{tabular}
\par\smallskip{\footnotesize Completa tiene siempre 20.652 coordenadas; las otras estrategias pueden elegir subconjuntos distintos entre particiones.}
\end{resulttable}
\FloatBarrier


GFS obtiene mayor macro-F1 medio que Individual en las ocho configuraciones, pero el método con mayor desempeño depende del dataset. Completa alcanza los mayores valores en CUReT y en Outex con SVM; GFS obtiene el mayor valor en Outex con ResMLP. En DTD destaca GFS, mientras que en FMD destacan Top-$k$ con SVM y GFS con ResMLP. Esta comparación describe las medias externas y no implica que cada diferencia sea estadísticamente significativa.

La Tabla~\ref{tab:dimensiones} permite contrastar el desempeño con el tamaño de la representación en las cinco estrategias. La dimensión media de GFS por configuración dataset--clasificador se sitúa entre 2.909,3 y 5.632 componentes, una reducción de entre 71,7\% y 85,4\% respecto de los 19.884 de Completa. Este rango compara las ocho medias, no los extremos de las particiones individuales. La disminución corresponde al vector almacenado y utilizado por el clasificador. No se midió una reducción proporcional de latencia o memoria de extremo a extremo.

\section{Rankings y comparaciones estadísticas}\label{sec:replica-paiva}

Las Tablas~\ref{tab:rankings-comparison} y~\ref{tab:global-tests-comparison} presentan los rankings y los contrastes de la biblioteca de 21 descriptores bajo las dos formas de agregación.

\begin{resulttable}
\centering\fontsize{8.5}{10.5}\selectfont\setlength{\tabcolsep}{3.2pt}
\caption{Rankings medios de macro-F1 con las dos formas de agregación. Menor ranking indica mejor posición relativa.}
\label{tab:rankings-comparison}
\begin{tabular}{@{}lrrrr@{}}
\toprule
Estrategia & Friedman: dataset & Quade: dataset & Friedman: dataset--clasificador & Quade: dataset--clasificador \\
\midrule
Individual & $4{,}750$ & $4{,}800$ & $4{,}625$ & $4{,}611$ \\
Completa & $2{,}750$ & $2{,}800$ & $2{,}875$ & $2{,}944$ \\
Top-$k$ & $2{,}000$ & $2{,}200$ & $2{,}125$ & $2{,}250$ \\
GFS & $2{,}250$ & $1{,}700$ & $2{,}125$ & $1{,}639$ \\
Homogénea & $3{,}250$ & $3{,}500$ & $3{,}250$ & $3{,}556$ \\
\bottomrule
\end{tabular}
\par\smallskip{\footnotesize Dataset: SVM y ResMLP se promedian dentro de cada dataset (cuatro bloques). Dataset--clasificador: se conservan separados (ocho bloques).}
\end{resulttable}
\FloatBarrier

\begin{resulttable}
\centering\fontsize{8.5}{10.5}\selectfont\setlength{\tabcolsep}{4pt}
\caption{Contrastes globales de macro-F1 bajo las dos formas de agregación.}
\label{tab:global-tests-comparison}
\begin{tabular}{@{}llrrr@{}}
\toprule
Agregación & Prueba & Estadístico & gl & $p$ \\
\midrule
Dataset & Friedman & 7,6000 & 4 & 0,1074 \\
Dataset & Quade & 2,8632 & 4,12 & 0,0705 \\
Dataset--clasificador & Friedman & 13,6000 & 4 & 0,0087 \\
Dataset--clasificador & Quade & 5,1248 & 4,28 & 0,0032 \\
\bottomrule
\end{tabular}
\par\smallskip{\footnotesize En la primera agregación se promedian SVM y ResMLP por dataset; en la segunda se mantienen como bloques separados.}
\end{resulttable}
\FloatBarrier

\begin{resulttable}
\centering\fontsize{8.5}{10.5}\selectfont\setlength{\tabcolsep}{4pt}
\caption{Comparaciones frente a Individual: p-values ajustados mediante Holm bajo las dos formas de agregación.}
\label{tab:vs-individual-comparison}
\begin{tabular}{@{}lrr@{}}
\toprule
Comparación & Dataset & Dataset--clasificador \\
\midrule
GFS vs. Individual & $0{,}2281$ & $0{,}0157$ \\
Top-$k$ vs. Individual & $0{,}1391$ & $0{,}0157$ \\
Completa vs. Individual & $0{,}5891$ & $0{,}2149$ \\
Homogénea vs. Individual & $1{,}0000$ & $0{,}5739$ \\
\bottomrule
\end{tabular}
\par\smallskip{\footnotesize Dataset utiliza cuatro bloques; dataset--clasificador utiliza ocho. En ambos casos se aplica Holm sobre la familia de diez comparaciones.}
\end{resulttable}
\FloatBarrier


La Tabla~\ref{tab:vs-individual-comparison} presenta las mismas cuatro comparaciones frente a Individual en paralelo. Al promediar SVM y ResMLP por dataset, los valores ajustados con Holm son 0,2281 para GFS y 0,1391 para Top-$k$; al mantener separados los clasificadores, ambos alcanzan 0,0157. Completa y Homogénea obtienen 0,2149 y 0,5739 en la segunda organización. Como magnitud descriptiva, las diferencias absolutas de macro-F1 de GFS frente a Individual, promediando ambos clasificadores, son 0,0277 en DTD, 0,0066 en FMD, 0,0024 en CUReT y 0,0392 en Outex.

Las dos organizaciones responden a preguntas distintas: la primera resume el comportamiento medio por dataset y la segunda permite observar si el efecto se mantiene por clasificador. Las ocho filas desagregadas reutilizan los mismos datasets e imágenes, por lo que sus resultados deben interpretarse junto con la estructura de dependencia del protocolo. Los procedimientos adicionales y sus salidas Java se conservan en los materiales de reproducibilidad.
